\documentclass{jsarticle}
\usepackage[dvipdfmx, final]{graphicx}
\usepackage{amsmath}
\usepackage{amssymb}
\usepackage{chemfig}
\begin{document}
\title{DNA of Universe and human  \\
being entrance with zeta function}
\author{}
\date{}
\maketitle
   \newcommand{\variint}{%
	\begin{picture}(15,30)
		\put(0,0){
			$ \iint $}
		\put(0,5){\line(15,0){15}}
	\end{picture} %
}

\newcommand{\variiint}{%
	\begin{picture}(15,15)
		\put(0,0){
			$ \iiint $}
		\put(0,5){\line(15,0){20}}
	\end{picture} %
}



 \newcommand{\alearletter}{%
	 \begin{picture}(10,20)
		 \put(0,0){
			 $ \square $}
		 \put(0,0){\line(1,1){10}}
	 \end{picture} %
	 }

 \newcommand{\secproduct}{%
	 \begin{picture}(10,20)
		 \put(0,0){
			 $ \nabla $}
		 \put(5,0){+}
	 \end{picture} %
	 }
 \newcommand{\talot}{%
	 \begin{picture}(15,20)
		 \put(0,0){
			 $ \text{\scalebox{2.0}[2]{$\nabla$}}$}
		 \put(5,3){
			 $ \Delta $}
	 \end{picture} %
	 }
 \newcommand{\hazerdnabla}{%
	 \begin{picture}(15,20)
		 \put(0,0){
			 $ \text{\scalebox{2.0}[2]{$\nabla$}}$}
		 \put(5,3){
			 $\text{Y}  $}
	 \end{picture} %
	 }

1998年度の富山大学入学試験の数学を、振り返って、
         第1問の問題　
	 対数関数の問題
         $$ y = \log x $$
	 第2問の問題
	 ベータ関数をグラフを書かせて、回転体の体積を求めさせて、
	 ゼータ関数と同型と導く問題
	 $$ \zeta(s) = {\beta(p,q) \over {\log x}} = 
	 {{\zeta(1-s)\zeta(0)\zeta(1)\zeta(2)\cdots } \over {(x_{1} - y_{1}) 
	 (x_{2} + y_{2})(x_{3} - y_{3})(x_{4} + y_{4})\cdots }} $$
	 第3問の問題
	 非線形の微分と積分の問題、
	 フェルマーの定理の融合問題、上の問題を融合させると、
	 一般相対性理論とAdS5多様体の式になる問題
	 $$ x^{n} + y^{n} -n xyz = 0, ||ds^2|| = e^{-2\pi T||\psi||}
	 [\eta + \bar{h}(x)]dx^{\nu}dx^{\mu} + T^2d^2 \psi $$
	 $$ x^n + y^n = [\eta + \bar{h}(x)],  z^n = T^2d^2 \psi $$

	 それらが、大域的微分と大域的積分多様体という分野を作る問題
	 になっている。
	 $$ \zeta(s) = {\beta(p,q) \over {\log x}}, \int \zeta(s) dx_m =
	 {\beta(p,q) \over {\log x}} $$
	 $$ = \int \zeta(s) e^{-f}dV $$
	 $$ e^{\pi} \cong \pi^{e}, \log x = {\beta(p,q) \over {\zeta(s)}} $$
	 $$ \zeta(s) = {\beta(p,q) \over {\log x}} $$
	 結局、ゼータ関数は、シャノンの公式であるという結論
	 $$ \zeta(s) = x\log x $$
	 $$ {\pi \over {e}} = \log \pi $$
	 $$ \pi = e \log \pi $$
	 ジョン・ナッシュさんが、提示した式を逆転させて、倍にすると、
	 ノルムのゼータ関数になる。
	 それは、シャノンの公式である。
         $$ ||ds^2|| = || x^{1 \over 2}\log x||^{4}  $$
	 $$ = x^{-{1 \over 2}} = || x^{1 \over 2} \log x||^{4} $$
	 $$ x^2 = ||x^2 \log x^{-1}||^{-4} $$
	 $$ x^2 = ||x^{-8} \log x^{4}|| $$
	 $$ ||ds^2|| = ||{x^{1 \over 2}  \log x^2}||^2 $$
	 $$ \zeta(s) = (x\log x)^n $$
	 これらが、サーストン・ペレルマン多様体になる。
	 $$ E(\sigma) = K(\sigma) \times H(\sigma) $$
	 $$ 4(\pi^{n}, e^{n}), x^n + y^n = z^n $$
	 これらが、一般相対性理論となるという結論
	 $$ \int \kappa T^{\mu\nu} \mathrm{dvol} = \int (R + {1 \over 2}
	 g_{ij}\Lambda)dx_m = \int{(R + {1 \over 2}g_{ij}\Lambda)}\mathrm{
		 dvol} $$
	 $$ e^{(R+{1\over 2}g_{ij}\Lambda)\log{( R+{1\over 2}g_{ij}
	 \Lambda)}} $$
	 $$ = e^{x\log x} = \int \Gamma^{'}(\gamma)dx_m = e^{-x\log x} $$
	 $$ x^n + y^n = e^{f} + e^{-f} \ge e^{f} - e^{-f} $$
	 $$ {d \over {d f}}F(x,y) = F^{f^{'}} = x^n + y^n $$
	 $$ =  e^{f} + e^{-f} \ge e^{f} - e^{-f} $$
	 $$ = \beta(p,q) = {\beta(p,q) \over {\log x}} = {\beta(p,q) \over 
	 \zeta(s)} $$

	 一般相対性理論の多様体積分は、宇宙のDNAであり、
	 ゼータ関数は、人体のDNAであるという結論、
	 ゼータ関数は、宇宙の広中先生の電話帳の住所録である。
	 人体の住所録でもある。
	 オイラーの定数の多様体積分が、それである。



	 $e^{\pi} $は、3次元球面、$\pi^{e}$は、3次元双曲面、
	 　$e+\pi$は、3次元射影平面、$e*\pi$は、3次元トーラス、
	 $e-\pi$
	 は、3次元クライン空間、${\pi \over e}$は、3次元Real Projective Plane,
	 ${e \over {\pi}}$、3次元Poincare球、$\sqrt{{e^{\pi}}}$は、3次元
	 Lens空間、$\sqrt{\pi^{e}}$は、3次元Seifert多様体、$e^{({\pi \over 2})
	 }$、3次元Brieskorn多様体であることと、
	 $$ ||ds^2|| = \pi r^2 = \pi {\left(||\int e^{\sin x\cos x} \int \sin x \cos x)dx ||\right)^{2}} = \pi {\left({x^{1 \over 2}\log x}^{4}\right)}^{2}  $$
	 $$ = \zeta(s) $$
	 レオナルド・オイラーのゼータ関数の$\pi$で表せられる、ゼータ関数の
	 サーストン・ペレルマン多様体の一般系の力学系である、
	 オイラーのゼータ関数の一般表示の方程式であるという、
	 レポートである。
\end{document}
